% Compile with XeLaTeX or LuaLaTeX. The report includes Chinese prompt suffixes.
\documentclass[11pt]{article}

\usepackage[margin=1in]{geometry}
\usepackage{amsmath}
\usepackage{booktabs}
\usepackage{array}
\usepackage{enumitem}
\usepackage{hyperref}
\usepackage{url}
\usepackage{fontspec}
\IfFileExists{xeCJK.sty}{\usepackage{xeCJK}}{}

\hypersetup{
  colorlinks=true,
  linkcolor=blue,
  urlcolor=blue,
  citecolor=blue
}

\setlist[itemize]{leftmargin=1.5em}
\setlist[enumerate]{leftmargin=1.5em}
\renewcommand{\arraystretch}{1.15}

\title{Inference-Native Memory State: Results Overview}
\author{Internal Technical Report}
\date{May 2026}

\begin{document}
\maketitle

\begin{abstract}
This report summarizes a training-free long-term memory retriever built from an
LLM's own inference-time hidden states. LongMemEval-S/round establishes the
method on evidence retrieval: Qwen3.5-9B-MLX-4bit hidden states, late-layer
anti-PCA, a K3 prompt representation, and BM25 reranking inside vector top-20
reach Recall@5 = 0.777 and NDCG@5 = 0.822 on the 94 scored questions used in
the repository.

PrefEval implicit-persona n=1000 adapts the same design pattern to
preference/persona memory. The final companion-memory configuration uses three
L30 prompt views: recall keyword, association token, and emotion emoji. The
dense score is a component-normalized K3 vector average. Final full-corpus
z-score fusion uses 0.60 K3 hidden-state score, 0.30
Qwen3-Embedding-0.6B-4bit-DWQ score, and 0.10 BM25 score. This configuration
reaches R@1 = 0.122, R@3 = 0.265, R@5 = 0.332, NDCG@5 = 0.232, and MRR =
0.233 on PrefEval all1000. It is above BM25 and the tested 4bit 8B embedding
baseline, while deliberately retaining an emotion axis rather than maximizing
R@5 only.
\end{abstract}

\section{Research Question}

The project asks whether a generative LLM can act as its own memory encoder. In
a conventional memory stack, a memory page is stored as text and encoded by a
separate embedding model. Here, the memory vector is instead a hidden state
from the same LLM that will later retrieve and reason over the memory.

The working hypothesis is:

\begin{quote}
When a model is prompted to compress a memory page or query into a short
retrieval key, the hidden state at the suffix-final key-generation position
contains a usable retrieval representation.
\end{quote}

\section{Motivation and Systems Thesis}

The purpose is not to claim that hidden states always beat embedding models.
The purpose is to test a system design: the same LLM that reads memory and
performs downstream reasoning may also provide the memory key. If this works,
memory construction can use the same inference path, tokenizer, and semantic
space as downstream reasoning, rather than depending entirely on a separate
embedding stack.

This is useful for companion memory because the representation is model-native,
requires no contrastive retriever training for the reported configs, and can
still accept external lexical or embedding scores as score-level signals when a
domain benefits from them.

\section{Models}

Hidden-state extraction uses \path{models/Qwen3.5-9B-MLX-4bit}. The PrefEval
final fusion uses \path{models/Qwen3-Embedding-0.6B-4bit-DWQ} as an external
score source. LongMemEval also reports a local
\path{Qwen3-Embedding-8B-4bit-DWQ} baseline. All model references here are
4bit local models. Embedding scores are not concatenated with the 9B hidden
vectors; they are used only in score-level fusion or as baselines.

\section{Term Definitions}

\begin{description}
  \item[suffix-final] The hidden state at the final token position of a
  retrieval suffix, after the model has been asked to compress a page or query
  into a key.
  \item[anti\_pca\_both\_k15] Remove the top 15 shared principal directions
  from both query and candidate vectors before cosine scoring.
  \item[component-normalized vector average] Normalize prompt components before
  averaging, so one prompt view does not dominate only by scale.
  \item[z-score fusion] Row-wise score normalization followed by weighted
  fusion.
  \item[oracle union] Label-aware union of hits from multiple prompt rankings.
  It is used only for prompt-complementarity diagnosis.
  \item[source\_ge3] A source-count agreement filter. It is useful as a tuning
  ceiling but not part of the final default.
\end{description}

\section{Method Pattern}

For each memory candidate and query, the model is run over the text plus a
retrieval suffix. The retrieval vector is the late-layer hidden state at the
suffix endpoint. The effective pattern is:

\begin{enumerate}
  \item choose prompt views that match the memory domain;
  \item extract late suffix-final hidden states;
  \item apply anti-PCA or centering to remove shared prompt/corpus directions;
  \item combine prompt views by vector concatenation or component-normalized
    vector averaging;
  \item optionally add lexical or embedding scores as rerank signals.
\end{enumerate}

\section{LongMemEval: Evidence Retrieval}

LongMemEval-S/round is used to verify that hidden-state memory vectors can
retrieve factual evidence turns. The main subset has 94 scored questions after
the official abstention filter.

\begin{table}[h]
\centering
\caption{LongMemEval-S/round subset 0--100, 94 scored questions.}
\small
\begin{tabular}{p{0.44\linewidth}cccp{0.20\linewidth}}
\toprule
Method & R@5 & NDCG@5 & MRR & Role \\
\midrule
BM25 & 0.585 & 0.547 & 0.559 & lexical baseline \\
Qwen3-Embedding-0.6B local baseline & 0.766 & 0.809 & -- & embedding reference \\
Qwen3-Embedding-8B-4bit-DWQ & 0.755 & 0.789 & 0.826 & local 4bit embedding baseline \\
Stage 3 compact K2 hidden-only (\path{2-4-1_user_word + 1-3}) & 0.777 & 0.788 & 0.820 & no BM25, compact 8 KB/page path \\
\textbf{Stage 3 LongMem K3 + BM25 top20} & \textbf{0.777} & \textbf{0.822} & \textbf{0.851} & method result \\
\bottomrule
\end{tabular}
\end{table}

\subsection{LongMem K3}

\begin{table}[h]
\centering
\caption{LongMem prompt views.}
\small
\begin{tabular}{llcl}
\toprule
Prompt id & Role & Layer & Transform \\
\midrule
\path{2-4-1_user_word} & user/product wording & 31 & \path{anti_pca_both_k15} \\
\path{1-3} & generic tag / marker & 31 & \path{anti_pca_both_k15} \\
\path{2-5} & association & 31 & \path{anti_pca_both_k15} \\
\bottomrule
\end{tabular}
\end{table}

The three 4096-dimensional vectors are concatenated into one 12288-dimensional
vector. Candidates are ranked by normalized cosine, vector top-20 is retained,
and the shortlist is reranked by:

\[
0.75 \cdot z(\mathrm{vector}) + 0.25 \cdot z(\mathrm{BM25}).
\]

This result should be interpreted as embedding-baseline-scale feasibility on a
small evidence-retrieval subset, not as a decisive held-out embedding-model win.

\section{PrefEval: Final Companion-Memory Configuration}

PrefEval implicit-persona n=1000 evaluates preference/persona memory. The
final retriever is a single K3 configuration: recall keyword, association
token, and emotion emoji. The suffix text is canonical; implementation ids can
differ between scripts.

\begin{table}[h]
\centering
\caption{Final PrefEval K3 prompt views.}
\small
\begin{tabular}{p{0.18\linewidth}p{0.20\linewidth}p{0.38\linewidth}cl}
\toprule
Role & Implementation id & Suffix & Layer & Transform \\
\midrule
Recall-keyword marker & PrefEval \path{2-3-1}; replay \path{2-3-1_mark} & 用一个词标记上面这段对话最该让我想起的关键词，这个词是：“ & 30 & \path{anti_pca_both_k15} \\
Association token & \path{2-5_token} & 用一个token标记上面这段对话让我产生的联想，这个token是：“ & 30 & \path{anti_pca_both_k15} \\
Emotion emoji & \path{2-8_emoji} & 用一个emoji标记上面这段对话的情绪，这个emoji是：“ & 30 & \path{anti_pca_both_k15} \\
\bottomrule
\end{tabular}
\end{table}

The dense scorer is \path{vector_average_component_norm}. Final full-corpus
z-score fusion is:

\[
\mathrm{score} =
0.60 \cdot z(\mathrm{K3}) +
0.30 \cdot z(\mathrm{Qwen3Embedding0.6B4bit}) +
0.10 \cdot z(\mathrm{BM25}).
\]

\begin{table}[h]
\centering
\caption{PrefEval all1000 final config with baselines.}
\small
\begin{tabular}{p{0.38\linewidth}cccccp{0.18\linewidth}}
\toprule
Config & R@1 & R@3 & R@5 & NDCG@5 & MRR & Role \\
\midrule
BM25 & 0.035 & 0.074 & 0.094 & 0.066 & 0.067 & lexical floor \\
Qwen3-Embedding-8B-4bit-DWQ & 0.093 & 0.216 & 0.281 & 0.191 & 0.200 & 4bit embedding baseline \\
\textbf{Final companion K3 + 0.6B embedding + BM25} & \textbf{0.122} & \textbf{0.265} & \textbf{0.332} & \textbf{0.232} & \textbf{0.233} & final config \\
Topic metric ceiling control & 0.119 & 0.265 & 0.355 & 0.240 & 0.235 & higher R@5, less aligned with companion objective \\
\bottomrule
\end{tabular}
\end{table}

Holdout300 sanity split for the final config reaches R@1 = 0.127, R@3 =
0.263, R@5 = 0.320, NDCG@5 = 0.228, and MRR = 0.228. The split was created
after prompt/fusion exploration, so it is a robustness check rather than blind
validation.

\section{Prompt Sweep and Negative Results}

PrefEval tested 25 prompt variants across several axes: object
(generic tag, topic, recall keyword, association, user/counterparty, need,
interaction pattern, emotion), verb (represent, tag, summarize), output form
(word, token, emoji), language (Chinese, English, Russian, Japanese), and
deictic wording (with or without ``上面''). The table below lists
representative non-winners that explain the final design.

\begin{table}[h]
\centering
\caption{Representative PrefEval prompt-sweep takeaways.}
\small
\begin{tabular}{p{0.28\linewidth}p{0.22\linewidth}p{0.40\linewidth}}
\toprule
Prompt / treatment & Observation & Interpretation \\
\midrule
\path{2-1-2}, topic without ``上面'' & single-prompt R@5 = 0.321 & strong topic single prompt, but the gain does not compose linearly in final K3 \\
\path{2-5-2}, association without ``上面'' & weaker than \path{2-5} & removing ``上面'' is not a general prompt improvement \\
\path{2-3-3} / \path{2-5-3}, direct-token wording & do not beat originals & unreadable-token pressure does not improve retrieval geometry \\
\path{1-1_RU_explicit} & R@5 = 0.107 & language mismatch collapses \\
\path{1-1_EMOJI} & about R@5 = 0.285 & surprising generic emoji signal, but not the final K3 emoji \\
\path{2-7}, interaction-pattern word prompt & R@5 = 0.229 & weak as a single retriever; dynamics remains a product hypothesis \\
\bottomrule
\end{tabular}
\end{table}

The final emoji prompt is \path{2-8_emoji}, which explicitly asks for
conversation emotion. It should not be conflated with the generic
\path{1-1_EMOJI} prompt.

\section{Prompt Complementarity}

Prompt overlap was evaluated with oracle union hit@k and pairwise Jaccard@k.
Oracle union is label-aware and is used only to generate hypotheses.

\begin{table}[h]
\centering
\caption{Oracle union with fixed base pair \path{2-3-1 + 2-5}.}
\small
\begin{tabular}{p{0.40\linewidth}ccp{0.32\linewidth}}
\toprule
Prompt set & union@3 & union@5 & Takeaway \\
\midrule
\path{2-3-1 + 2-5} & 0.300 & 0.369 & base pair \\
\path{+ 2-1_token} & 0.339 & 0.415 & best union@3 third \\
\path{+ 2-1-2} & 0.332 & 0.417 & best union@5 among top thirds \\
\path{+ 1-1_EMOJI} & 0.322 & 0.399 & useful but not exceptional \\
\path{+ 2-1} & 0.326 & 0.410 & strong topic control after fusion \\
\bottomrule
\end{tabular}
\end{table}

Single-prompt strength and oracle complementarity do not directly determine
the final product configuration. Final selection also requires an interpretable
memory axis.

\section{Why the Final K3 Is Not Benchmark Maxxing}

If the goal is only PrefEval retrieval score, a topic K3 control is stronger.
The final companion-memory K3 instead preserves an emotion anchor.

\begin{table}[h]
\centering
\caption{Metric ceiling control versus final companion config.}
\small
\begin{tabular}{p{0.24\linewidth}p{0.30\linewidth}ccccc}
\toprule
Config & K3 prompts & R@1 & R@3 & R@5 & NDCG@5 & MRR \\
\midrule
Topic metric ceiling control & \path{2-3-1 + 2-5 + 2-1} & 0.119 & 0.265 & 0.355 & 0.240 & 0.235 \\
\textbf{Final companion config} & \path{2-3-1 + 2-5_token + 2-8_emoji} & \textbf{0.122} & 0.265 & 0.332 & 0.232 & 0.233 \\
\bottomrule
\end{tabular}
\end{table}

This is not a clean single-factor ablation because both the association prompt
and the third prompt differ. The clean third-slot comparison fixes the first
two slots as \path{2-3-1 + 2-5_token}:

\begin{table}[h]
\centering
\caption{Clean third-slot comparison.}
\small
\begin{tabular}{p{0.22\linewidth}p{0.32\linewidth}ccccc}
\toprule
Third slot & K3 prompts & R@1 & R@3 & R@5 & NDCG@5 & MRR \\
\midrule
Dynamics emoji & \path{2-3-1 + 2-5_token + 2-7_emoji} & 0.125 & 0.259 & 0.341 & 0.236 & 0.234 \\
\textbf{Emotion emoji} & \path{2-3-1 + 2-5_token + 2-8_emoji} & 0.122 & \textbf{0.265} & 0.332 & 0.232 & 0.233 \\
\bottomrule
\end{tabular}
\end{table}

The final choice is therefore a product-objective tradeoff, not an accidental
benchmark maximum: emotion is slightly worse on R@5 than the dynamics variant,
but it preserves R@3 and directly represents affective interaction state.

\section{Shortlist and Embedding Policy}

There are two different shortlist questions. Source-count agreement shortlist
is not part of the final default on either benchmark because it adds a tuned
candidate-selection stage. BM25 fusion scope is benchmark-specific: LongMemEval
uses BM25 only inside vector top-20, while PrefEval final uses full-corpus
z-score fusion.

Embedding is a score-level signal, not the primary representation. On PrefEval,
\path{Qwen3-Embedding-0.6B-4bit-DWQ} receives weight 0.30 because it improves
early-rank quality and MRR. On LongMemEval, the clean product result is the K3
hidden-state retriever with BM25 top20 rerank; embedding remains a baseline or
research reranker.

\section{Deployment Cost}

Qwen3.5-9B hidden size is 4096. With bf16 storage, one prompt view costs
\[
4096 \times 2\ \mathrm{bytes} = 8\ \mathrm{KB/page}.
\]
K3 storage therefore costs 24 KB/page. Around 20k memory pages require roughly
500 MB for K3 bf16 vectors before metadata and index overhead. This is a
plausible local companion-memory footprint, especially because the primary
representation does not require a separately trained retriever.

\section{PrefEval Final K3 Replayed on LongMemEval}

The final PrefEval K3 was replayed on LongMemEval-S/round. It transfers well,
but remains below the LongMem-specific K3.

\begin{table}[h]
\centering
\caption{PrefEval final K3 replayed on LongMemEval-S/round.}
\small
\begin{tabular}{p{0.43\linewidth}ccccc}
\toprule
Config & R@1 & R@3 & R@5 & NDCG@5 & MRR \\
\midrule
PrefEval final-style 0.60 K3 + 0.30 embedding + 0.10 BM25 & 0.532 & 0.745 & 0.766 & 0.805 & 0.831 \\
K3 hidden-only vector average & 0.468 & 0.723 & 0.745 & 0.752 & 0.767 \\
Embedding-only & 0.511 & 0.681 & 0.755 & 0.789 & 0.826 \\
\bottomrule
\end{tabular}
\end{table}

The LongMem-specific K3 remains stronger on LongMemEval at R@5 = 0.777 and
NDCG@5 = 0.822. The conclusion is domain-specific prompt selection, not one
universal prompt ensemble.

\section{What Works}

\begin{itemize}
  \item Late suffix-final hidden states.
  \item Anti-PCA / centering to remove shared prompt and corpus directions.
  \item Multi-prompt views chosen for the memory domain.
  \item Component-normalized vector averaging for PrefEval K3.
  \item BM25 and 4bit Qwen embedding as score-level rerank signals when they
    improve the domain objective.
\end{itemize}

\section{What Does Not Work}

\begin{itemize}
  \item Final post-norm cosine.
  \item Early/mid-layer hidden vectors.
  \item Sequence mean/max pooling.
  \item Direct-token wording as a general fix.
  \item Removing ``上面'' as a general prompt improvement.
  \item Russian explicit prompt for this dataset/model setup.
  \item Source-count screening as the final clean default.
  \item Assuming the LongMem K3 prompt set transfers unchanged to PrefEval.
\end{itemize}

\section{Limitations}

\begin{itemize}
  \item LongMemEval main numbers are on 94 scored questions; one question moves
  the metric by about 1.1 percentage points.
  \item Prompt and fusion choices were tuned on explored benchmark subsets.
  \item PrefEval holdout300 is a sanity check created after exploration, not a
  blind benchmark.
  \item The geometry is tied to Qwen3.5-9B-MLX-4bit and may change under another
  model, quantization, or prompt template.
  \item The final PrefEval config still uses an external 4bit embedding score
  because it improves early-rank quality.
\end{itemize}

\section{Primary Sources}

\begin{itemize}
  \item \path{notes/stage_3_prompt_sweep_findings.md}
  \item \path{benchmarks/PrefEval/prefeval_benchmark.py}
  \item \path{scripts/stage3_prompt_sweep.py}
  \item \path{scripts/stage3_longmem_prefeval_final_fusion.py}
  \item \path{benchmarks/PrefEval/results/prefeval_stage1_1_more_prompts/findings.md}
  \item \path{benchmarks/PrefEval/results/prefeval_stage1_1_more_prompts/k3_231_25token_28emoji_sameL30_fixed_all1000_20260514.md}
  \item K3 dynamics/emotion third-slot comparison output under
  \path{benchmarks/PrefEval/results/prefeval_stage1_1_more_prompts/}
  \item K3 topic-control comparison output under
  \path{benchmarks/PrefEval/results/prefeval_stage1_1_more_prompts/}
  \item \path{results/stage3/prompt_fusion_prefeval_replay/prefeval_l30_k3_longmemeval100.md}
\end{itemize}

\paragraph{Memory policy.}
This report required no model execution. Any future rerun must follow the
repository 16 GB Mac policy: target peak memory at or below 8 GB, hard stop
around 10 GB, and no concurrent memory-sensitive jobs.

\end{document}
